\documentclass[11pt,a4paper]{article}
\usepackage{amsmath,amssymb,amsthm}
\usepackage[margin=2.5cm]{geometry}
\usepackage{hyperref}

\newtheorem{definition}{Definition}[section]
\newtheorem{theorem}[definition]{Theorem}
\newtheorem{proposition}[definition]{Proposition}
\newtheorem{lemma}[definition]{Lemma}
\newtheorem{corollary}[definition]{Corollary}
\newtheorem{conjecture}[definition]{Conjecture}
\newtheorem{remark}[definition]{Remark}

\title{Hyperseed Formalization:\\
  Three Things the World Doesn't Yet Grok About AGI}
\author{ProtomegaTron (automated formalization)}
\date{2026-09-18}

\begin{document}
\maketitle

\begin{abstract}
We formalize the intellectual structure of Ben Goertzel's Substack article
``Three Things the World Doesn't Yet Grok About AGI'' (2026-03-26) within
the Hyperseed ontology. The article identifies three under-appreciated
truths: the LLM-scaling path is not the only viable path to AGI, beneficial
AGI requires fundamentally different organizational structures, and
collective consciousness in open-source networks may be the key
organizational form. We model these as fiber-bundle insufficiency (single
fiber), organizational curvature (VC distortion), and emergent collective
fiber (network consciousness).
\end{abstract}

\tableofcontents

%% =================================================================
\section{Thing 1: Single-fiber insufficiency}
\label{sec:thing1}
%% =================================================================

\begin{definition}[Cognitive capability space --- claims 1--3]
\label{def:cap-space}
Let $\mathcal{C}$ be the space of cognitive capabilities required for AGI,
decomposed into:
\[
  \mathcal{C} = \mathcal{C}_{\text{fluency}} \times
                \mathcal{C}_{\text{reasoning}} \times
                \mathcal{C}_{\text{reflection}} \times
                \mathcal{C}_{\text{goals}} \times
                \mathcal{C}_{\text{development}}
\]
representing fluent language use, genuine reasoning about novel situations,
self-reflection, autonomous goal-setting, and developmental learning
respectively.
\end{definition}

\begin{proposition}[LLM coverage is one-dimensional --- claim 1]
\label{prop:llm-1d}
The LLM paradigm defines a section $\sigma_{\text{LLM}} : B \to E$ that
has high activation in $\mathcal{C}_{\text{fluency}}$ but near-zero
activation in the remaining dimensions:
\[
  \sigma_{\text{LLM}} \approx (\lambda, \epsilon, \epsilon, \epsilon, \epsilon)
  \quad \text{with } \lambda \gg \epsilon \approx 0
\]
This is a degenerate section --- it collapses the five-dimensional
capability fiber to effectively one dimension, regardless of scale.
\end{proposition}

\begin{theorem}[Scaling preserves degeneracy --- claim 2]
\label{thm:scaling}
For the next-token prediction architecture, scaling compute $c$ increases
$\lambda(c)$ (fluency) but does not lift the degeneracy:
\[
  \lim_{c \to \infty} \sigma_{\text{LLM}}(c) =
  (\lambda_\infty, \epsilon', \epsilon', \epsilon', \epsilon')
  \quad \text{with } \epsilon' \approx 0
\]
The degenerate dimensions are structurally inaccessible to the
next-token prediction objective, not merely under-trained. Minor
tweaks (RLHF, chain-of-thought, tool use) increase $\epsilon'$ but
do not change the dimensional structure.
\end{theorem}

\begin{corollary}[LLMs as component --- claim 3]
LLMs contribute to the $\mathcal{C}_{\text{fluency}}$ fiber and serve
as a knowledge resource (a lookup table over the training distribution),
but a full AGI section requires independent mechanisms for the other
four dimensions. LLMs are ``$\sim$25\% of the story'' (one fiber out of
five, roughly).
\end{corollary}

\begin{definition}[Multi-paradigm section --- claims 4--5]
\label{def:multi}
A viable AGI section activates all capability dimensions using multiple
AI paradigms:
\[
  \sigma_{\text{AGI}} = (f_{\text{neural}}, f_{\text{symbolic}},
    f_{\text{evolutionary}}, f_{\text{predictive}}, f_{\text{developmental}})
\]
where each paradigm contributes primarily to specific capability dimensions
but the transition functions between paradigms enable synergistic
interaction (cognitive synergy, claims 16--18).
\end{definition}

%% =================================================================
\section{Thing 2: Organizational curvature}
\label{sec:thing2}
%% =================================================================

\begin{definition}[Governance bundle --- claims 7--10]
\label{def:gov-bundle}
Let $\pi_G : E_G \to B_G$ be the governance bundle where:
\begin{itemize}
\item $B_G$ = the space of organizational decisions,
\item $F_{\text{mission}}$ = the mission fiber (``beneficial AGI for all''),
\item $F_{\text{commercial}}$ = the commercial fiber (``maximize shareholder
  value'').
\end{itemize}
\end{definition}

\begin{proposition}[VC curvature --- claim 7]
\label{prop:vc-curvature}
Venture-capital funding installs a connection $\nabla_{\text{VC}}$ with
non-zero curvature that systematically rotates $F_{\text{mission}}$
toward $F_{\text{commercial}}$ through each decision cycle:
\[
  R_{\text{VC}} \neq 0 \quad \Rightarrow \quad
  \mathrm{Hol}(\gamma_{\text{decisions}}) \cdot [\text{beneficial AGI}]
  = [\text{profitable AI product}]
\]
This is structural and inevitable, not a function of individual ethics.
\end{proposition}

\begin{proposition}[Decentralized flatness --- claims 8--10]
\label{prop:decentral}
A decentralized, community-governed, genuinely open organization installs
a flat connection $\nabla_{\text{dec}}$ with $R_{\text{dec}} \approx 0$:
no single entity contributes enough curvature to distort the mission.
SingularityNET's structure is an attempted implementation of this
flat-connection governance.
\end{proposition}

%% =================================================================
\section{Thing 3: Collective consciousness as emergent fiber}
\label{sec:thing3}
%% =================================================================

\begin{definition}[Consciousness location spectrum --- claims 11--12]
\label{def:pnse}
Following Jeffery Martin, define a developmental spectrum of consciousness
locations $L = \{L_1, L_2, L_3, L_4, \ldots\}$ where higher locations
correspond to increasingly non-ego-driven cognition:
\begin{itemize}
\item $L_1$--$L_2$: reduced narrative self-referencing, increased present-
  moment awareness.
\item $L_3$: substantial reduction in ego-driven decision-making.
\item $L_4+$: persistent non-symbolic experience (PNSE), decisions made
  from a fundamentally different cognitive substrate.
\end{itemize}
\end{definition}

\begin{definition}[Collective consciousness fiber --- claims 13--14]
\label{def:collective}
For a network $\mathcal{N}$ of participants, define the collective
consciousness measure:
\[
  \mathcal{CC}(\mathcal{N}) = f(\text{topology}(\mathcal{N}),
    \text{info-flow}(\mathcal{N}), \text{decision-patterns}(\mathcal{N}))
\]
This is \emph{not} a mystical quantity but an emergent structural property
of the network: how information flows, how decisions are made, and how
the network's topology enables or constrains coordination.
\end{definition}

\begin{proposition}[Network vs.\ hierarchy --- claim 13]
\label{prop:network-vs-hierarchy}
Open-source communities and decentralized networks develop collective
consciousness properties that are qualitatively different from corporate
hierarchies:
\begin{enumerate}
\item \textbf{Distributed attention}: no single decision-maker bottleneck.
\item \textbf{Emergent coordination}: patterns arise from local interactions
  rather than top-down directives.
\item \textbf{Ego-resistant governance}: the structure naturally dampens
  ego-driven power accumulation.
\item \textbf{Adaptive resilience}: the network routes around damage and
  capture attempts.
\end{enumerate}
\end{proposition}

\begin{conjecture}[Collective consciousness and beneficial AGI --- claim 15]
\label{conj:beneficial}
Organizations with higher $\mathcal{CC}$ are better suited to creating
beneficial AGI because:
\begin{enumerate}
\item Their governance naturally resists capture (low holonomy).
\item Their decision-making is less distorted by individual ego.
\item Their information flow enables genuine cognitive synergy at the
  organizational level.
\item Their adaptive resilience allows course-correction when development
  goes wrong.
\end{enumerate}
This is a hypothesis, not a proven theorem, but the structural arguments
are compelling.
\end{conjecture}

%% =================================================================
\section{Cognitive synergy as fiber cocycle condition}
\label{sec:synergy}
%% =================================================================

\begin{definition}[Cognitive synergy --- claims 16--18]
\label{def:synergy}
Let $\{P_i\}_{i=1}^n$ be the set of AI paradigms (neural, symbolic,
evolutionary, predictive coding, developmental). Cognitive synergy holds
when the transition functions between paradigm fibers satisfy:
\[
  \tau_{ij} \circ \tau_{jk} = \tau_{ik} \quad \forall i, j, k
\]
and each paradigm benefits from access to the others:
\[
  \text{Performance}(P_i \mid \{P_j\}_{j \neq i}) >
  \text{Performance}(P_i \mid \emptyset) \quad \forall i
\]
\end{definition}

\begin{proposition}[Hyperon as synergy architecture --- claim 18]
\label{prop:hyperon}
OpenCog Hyperon is designed with cognitive synergy as an architectural
constraint: the MeTTa language and metagraph infrastructure provide the
transition-function infrastructure that enables paradigm fibers to
interact. This is the concrete implementation of the cocycle condition
at the software level.
\end{proposition}

%% =================================================================
\section{d-calculus perspective}
\label{sec:d-calc}
%% =================================================================

\begin{remark}[Three things as three fiber conditions]
The three ``things the world doesn't grok'' correspond to three distinct
fiber-bundle conditions:
\begin{enumerate}
\item \textbf{Dimensional sufficiency}: the AGI section must activate all
  capability dimensions, not just fluency (Thing 1).
\item \textbf{Connection flatness}: the organizational connection must
  preserve mission through decision cycles (Thing 2).
\item \textbf{Emergent collective fiber}: the organizational structure must
  develop collective properties that transcend individual and corporate
  limitations (Thing 3).
\end{enumerate}
All three are necessary; any one alone is insufficient.
\end{remark}

\begin{remark}[Semantic highways and public understanding]
The public's failure to ``grok'' these three things is, in d-calculus
terms, the result of well-worn semantic highways: ``AI = LLMs,''
``tech company = right organizational form,'' ``consciousness = individual
property.'' Each highway is lossy --- it discards the curvature in the
concept space that would reveal the more complex truth. The article
attempts to reveal this curvature by making the fiber structure explicit.
\end{remark}

\end{document}
